\inicializarSubseccion{Fórmulas, ecuaciones y constantes}
Constante de Coulomb:
\begin{equation}
    k = \frac{1}{4\pi\epsilon_0} \approx 8.99 \times 10^9 \, \text{N m}^2/\text{C}^2
\end{equation}

Distancia de un punto a cada carga:
\begin{equation}
    r_{\text{pos}} = \sqrt{(x - \frac{d}{2})^2 + y^2}
\end{equation}
\begin{equation}
    r_{\text{neg}} = \sqrt{(x + \frac{d}{2})^2 + y^2}
\end{equation}

Campo eléctrico de una carga puntual:
\begin{equation}
    E_{x\text{-pos/neg}} = \frac{k(\pm q)(x-x_{\text{pos/neg}})}{r_{\text{pos/neg}}^3}
\end{equation}
\begin{equation}
    E_{y\text{-pos/neg}} = \frac{k(\pm q)y}{r_{\text{pos/neg}}^3}
\end{equation}

Superposición vectorial:
\begin{equation}
    E_x = E_{x\text{-pos}} + E_{x\text{-neg}}
\end{equation}
\begin{equation}
    E_y = E_{y\text{-pos}} + E_{y\text{-neg}}
\end{equation}

Magnitud del campo eléctrico para visualización:
\begin{equation}
    E_{\text{mag}} = \sqrt{E_x^2 + E_y^2}
\end{equation}
\begin{equation}
    E_{x\text{-mag}} = \frac{E_x}{E_{\text{mag}}}
\end{equation}
\begin{equation}
    E_{y\text{-mag}} = \frac{E_y}{E_{\text{mag}}}
\end{equation}
